% ======================================================================
% Title, abstract, introduction.
% Narrative spine: an autonomous agent discovered new plausibility rules
% for inorganic crystal structures. Everything else in the paper is
% either what the rules are, or what the agent learned about how such
% rules can and cannot be discovered.
% ======================================================================

% ---------------------------------------------------------------- TITLE
% Candidates considered:
%   1. Autonomous discovery of calibrated plausibility rules for
%      inorganic crystal structures                            [CHOSEN]
%   2. An autonomous agent discovers crystal-structure plausibility
%      rules that never abstain
%   3. Machine discovery of contact-ratio rules for inorganic crystal
%      plausibility
%   4. Autonomous discovery, and self-refutation, of plausibility rules
%      for inorganic crystals
%   5. Agent-discovered plausibility rules for inorganic crystals, and
%      the counterexample bound on discovering them
% Chosen 1: names the object (inorganic crystal structures), the
% capability (autonomous discovery) and the differentiator (calibrated),
% is searchable, and every word is supported by the data.

\title[Autonomous discovery of crystal plausibility rules]{Autonomous discovery of
calibrated plausibility rules for inorganic crystal structures}

% ------------------------------------------------------------- ABSTRACT
\abstract{Rules that judge whether a proposed crystal structure is chemically plausible
are used to teach solid-state chemistry, to curate structural databases, and now to
filter the output of generative models for materials. The canonical set, formulated by
Pauling in 1929 from a few hundred structures, is violated by most oxides when its rules
are applied jointly, yet nothing has displaced it. Here we report plausibility rules
discovered autonomously by a machine agent from 99{,}162 experimental structures under a
pre-registered protocol with a sealed evaluation split. The leading rule bounds the
shortest cation--anion contact by the sum of the two ionic radii. It is satisfied by
99.2\% of held-out experimental crystals, is element- and valence-aware, costs one
distance computation and one table lookup, and, unlike Pauling's rules, never declines to
distinguish two candidate structures. An ionicity-conditioned upper bound repairs the one
corruption a lower bound cannot detect, and transfers to a corruption class the search
never saw. Enumerating provably optimal decision trees shows that the constraint on such
rules is not the space of expressible hypotheses but the counterexamples used to test
them. The agent also overturned eleven of its own conclusions; we publish that ledger
with the rules, because each entry leaves a diagnostic that outlives the claim it
killed.}

% --------------------------------------------------------- INTRODUCTION
\section{Introduction}

Deciding whether a proposed arrangement of atoms could be a real crystal is a judgement
that solid-state chemistry has to make thousands of times before it makes a single
measurement. It is made when a crystallographer assesses a refinement, when a database
curator screens a deposition, and, increasingly, when a generative model proposes a
structure that no one has yet attempted to make. The judgement is old enough to have a
canonical answer. Pauling's five rules, set down in 1929 from the few hundred structures
then available, ask whether cation--anion distances match a radius ratio, whether bond
strengths sum correctly at each anion, whether coordination polyhedra share too many
edges and faces, whether highly charged cations avoid one another, and whether a
structure uses as few distinct environments as it can \cite{pauling1929principles}. Each
reduces a crystal to arithmetic on a radius table and a set of formal charges, and each
returns a verdict in a minute with a pencil. Their tenure rests on that economy.

Economy is not accuracy, and the gap between them has now been measured. A 2020 audit of
inorganic structures found that only a small minority satisfy Pauling's second to fifth
rules simultaneously, that individual rules fare little better than chance in places, and
that the verdicts move substantially with the algorithm chosen to assign neighbours; its
authors closed by calling for machine-learned replacements \cite{george2020limited}. Six
years later the call is open. The obstacle is not a shortage of structural data or of
learning algorithms, both of which have grown by orders of magnitude
\cite{jain2013commentary,shannon1976revised}. It is that nobody has agreed what a
replacement should be measured against.

That difficulty is sharper than it first appears, because the metric the audit itself
uses cannot rank rules. A rule satisfied by almost every real crystal is trivially
obtainable: the statement that every interatomic distance is positive scores a perfect
satisfaction rate and excludes nothing. The consequences are not hypothetical. Generative
models for crystal structures are routinely validated by asking whether any pair of atoms
lies closer than a fixed absolute distance, a criterion that on our test set accepts every
real crystal while rejecting only a small percentage of deliberately corrupted ones. A
field-wide quality filter can therefore be very close to vacuous while reporting a
satisfaction rate near unity. Any rule, learned or hand-made, must be scored on how often
it accepts real crystals and on how often it rejects implausible ones, jointly.

Here we report plausibility rules discovered by an autonomous agent rather than proposed
by hand. Given 99{,}162 experimental structures, a corpus of derived structural
quantities, and a pre-registered protocol that sealed part of the data before any search
began, the agent proposed candidate predicates, evaluated them against both real and
deliberately corrupted structures, and kept those that survived. The rule it returns
first is a single inequality on a single dimensionless quantity, the ratio of the
shortest cation--anion distance in a structure to the sum of the two ionic radii
\cite{shannon1976revised}. Using a radius sum rather than the radius ratio of Pauling's
first rule is the substantive change: a ratio compares two lengths that never meet,
whereas a sum is the length that a contact is attempting to reach. The resulting
criterion is satisfied by 99.2\% of held-out experimental crystals, is stable to within a
fifth of a percentage point across three independent neighbour-assignment algorithms, and
commits on every structure pair we present to it. Extending it with an upper bound that
applies only where the ionic model itself applies repairs the single corruption mode a
lower bound cannot detect, and does so on a corruption class the selection procedure was
never shown.

What the agent learned about discovering such rules is the second result, and it is the
more transferable one. Enumerating the provably optimal decision tree at each point of a
cost trade-off, over the same features, yields rule sets that dominate the hand-built
ones in distribution, accepting more real crystals and rejecting more corrupted ones at
once. Withholding a single class of corruption removes that advantage entirely. The
binding constraint on a learned structural rule is therefore not the space of hypotheses
the search can express, which certified optimality lets us bound, but the counterexamples
used to test it: a rule trained against a synthetic corruption process learns the
corruption process, and the reward is indistinguishable from chemistry until a class is
withheld.

Autonomy here means a specific and checkable thing. The agent held its own evaluation
split sealed throughout, wrote its hypotheses down before testing them, and was required
to pair every claim with a control designed to kill it. That discipline is what separates
the present work from an automated parameter sweep, and it is also what produced the
paper's least comfortable result: the loop generates candidate conclusions far faster
than it generates the controls that refute them, so the scarce resource in autonomous
discovery is not search but the arithmetic of what a number is allowed to mean.

The same protocol that produced the rules also produced eleven conclusions that did not
survive, each overturned by the agent after it had been stated. We report them as a
ledger beside the surviving results. They are not a confession; each entry cost one cheap
diagnostic to expose, and those diagnostics apply well beyond the claims that motivated
them.
